% ===========================================================
% literature_review.tex
% ===========================================================
\section{Background and Related Work}
\label{sec:related}

\subsection{Locational marginal pricing}
The locational marginal price (LMP) construct on which this work
depends derives from the spot-pricing framework of
Schweppe~\textit{et al.}~\cite{Schweppe1988} and the contract-network
formulation of Hogan~\cite{Hogan1992}, who decomposed
\begin{equation}
\lambda_b = \lambda_{\mathrm{ref}} +
            \mathrm{LMC}_b + \mathrm{LCC}_b
\label{eq:lmp-decomp}
\end{equation}
into reference energy, marginal-loss component (LMC), and
marginal-congestion component (LCC).  The present pipeline absorbs
$\mathrm{LMC}_b$ into a per-bus loss factor and $\mathrm{LCC}_b$ into
the island partition itself.

\subsection{Marginal versus average emission factors}
Hawkes~\cite{Hawkes2010} demonstrated that marginal CO$_2$ factors
can differ from system averages by an order of magnitude and that
this gap materially affects energy-efficiency build/no-build
decisions.  Siler-Evans~\textit{et al.}~\cite{SilerEvans2012}
regressed historical EPA generator-level data to fit MEFs at the
Reliability Council level and showed strong diurnal/seasonal cycles.
Callaway~\textit{et al.}~\cite{Callaway2018} formalised the
distinction between operational and decision-relevant emission
factors; the present paper sits firmly in the operational/short-run
camp, computing factors against the realised CEN dispatch.

\subsection{ISO production tools}
PJM~\cite{PJMLME2024,PJMPrimer} publishes five-minute per-node
emission rates with the same energy/loss/congestion structure as the
LMP itself.  CAISO~\cite{CAISOOutlook} surfaces a system-level
emission curve at five-minute resolution but is not nodal.
WattTime~\cite{WattTime2024} ships a balancing-area Marginal
Operating Emission Rate with a 72-hour forecast but is neither nodal
nor unit-attributed.  NREL's Cambium~\cite{NRELCambium} publishes
\emph{long-run} factors for capacity-expansion studies, not
operational dispatch.  Electricity
Maps~\cite{ElectricityMaps2024,ElectricityMapsCritique} publishes
hourly per-zone consumption-based factors and discontinued its
marginal feed in 2024.

\subsection{Counterfactual marginals}
A research-grade refinement is the counterfactual evaluation of
Valenzuela~\textit{et al.}~\cite{ValenzuelaASL2022}, who exploit
implicit differentiation of the QP/LP optimality conditions to
recover an exact dynamic locational marginal emission factor from a
solved dispatch.  The present pipeline is empirical rather than
counterfactual, but its open-source structure with the PLEXOS
\texttt{.accdb} oracle is designed so that the same input set can
be fed into the \texttt{gtopt} SDDP solver to produce a
counterfactual companion in future work.

\subsection{Prior CEN-specific work}
For Chile, Bonvallet~\textit{et al.}~\cite{Bonvallet2016}
co-optimised hydrothermal scheduling under non-convex irrigation
constraints; Moreno~\textit{et al.}~\cite{Moreno2014} treated
distributed-generation deployment under uncertainty.  The
Coordinador's annual \textit{Reporte Art.~72}~\cite{CENArt722024}
aggregates emissions at the system level, and the 4eChile
guide~\cite{FourEChileGuide} documents a manual marginal-emission
estimation tool for CDM projects.  None publishes a per-bus
per-quarter signal.

\subsection{Methodological choice points}
Three choices dominate the design space.  The
\emph{merit-stack picker}: should the marginal unit be the
last-dispatched cleared at $\lambda_{\mathrm{ref}}$
(``maximum CV~$\le \lambda$''~\cite{SilerEvans2012}), or the
absolute-closest CV?  The pipeline picks the absolute-closest CV
because the closest-above is sometimes nearer than the
closest-below at low-CMG hours.  The \emph{forced-unit
treatment}: units flagged with \texttt{consigna}$\in$\{CI, MT,
PMT, PP, PS, PC\} are price-takers when at technical minimum but
remain valid candidates when interior.  The \emph{loss-factor
decomposition}: the pipeline uses CEN's published MLF (extracted
from the PID \texttt{.accdb}) when available, falling back to the
empirical $\lambda_b/\lambda_{\mathrm{ref}}$ ratio when it is not.
